\section{Related Work}\label{sec:rw}

\subsection{Caching and memoization libraries on the \jvm}

The widely adopted \jvm caching libraries---Guava's \code{CacheBuilder}~\cite{guava} and Caffeine~\cite{caffeine}---offer bounded caches, \ttl expiry, statistics, and weak-reference support. Both are consumed through an explicit builder API that materialises a cache object at every call site, and both carry dependency footprints sized for server workloads. Caffeine further implements W-TinyLFU~\cite{einziger2017tinylfu}, a state-of-the-art admission policy that outperforms pure \lru and \lfu on a broad range of workloads. Neither library is routinely used inside Android applications, because their builder-centric API is incompatible with the annotation-driven style favoured in the Android ecosystem and their dependency footprints conflict with Android's 64K DEX method limit.

Spring Cache~\cite{spring_cache} offers the annotation ergonomics (\code{@Cacheable}, \code{@CacheEvict}) that a refactoring tool would want, but it is inseparable from the Spring application context and has no presence on Android. Project Reactor~\cite{reactor_cache} adds a \code{.cache()} operator for reactive streams, but this solves a different problem (caching an emitted sequence) rather than caching a pure function's return value under an argument tuple. Function-level wrapper libraries such as~\cite{javamemoization} avoid the builder overhead by exposing a \code{Function<A, R>} interface, but require every caller to hold a functional reference rather than call the method directly---effectively replacing one refactoring with another. MemoizeLib's annotation-plus-bytecode-rewrite approach occupies the intermediate point in this design space: it requires neither a Spring-style container nor a call-site rewrite, which is the property we argue in Section~\ref{sec:intro} is necessary for an automated refactoring target.

\subsection{Build-time bytecode transformation for cross-cutting concerns}

Rewriting bytecode at build time to implement a cross-cutting concern has a substantial history in the \jvm community. AspectJ~\cite{aspectj} pioneered compile-time and load-time weaving for cross-cutting concerns such as logging, transactions, and security; MemoizeLib's class visitor pipeline can be viewed as a purpose-built aspect weaver for one aspect. On Android specifically, the \agp instrumentation API~\cite{agp_instrumentation} exposes \asm class visitors as first-class build steps, which is the hook through which MemoizeLib integrates. Prior work by Hilton et al.~\cite{hilton_bytecode} uses the same hook for logging and tracing.

Hilt~\cite{hilt} and Dagger~\cite{dagger} both use compile-time code generation, through \ksp or annotation processors, to materialise dependency graphs. Their design trade-off is the inverse of ours: they generate new classes at compile time, whereas MemoizeLib rewrites existing classes at build time. The compile-time-generation approach exposes the generated code to debuggers and other tooling; the build-time-rewrite approach used here makes possible the property that enabling memoization on an existing method requires no change whatsoever to its call sites or to its enclosing type---a property that is critical for the automated-refactoring use case that motivates this work.

\subsection{Performance engineering on Android}

Empirical studies of Android performance have identified specific APIs and code patterns that incur disproportionate time and energy costs~\cite{patterns,reduce_bytecode_energy,energy_smells_palomba,marco_refactorings_wild}. This literature has consistently established that software-level optimisations---library choice, data-structure choice, algorithm choice---dominate the available performance headroom on mobile hardware, and that developers routinely miss opportunities for local, mechanical refactorings~\cite{rodrigo_fix_static_warnings2019,empirical_rua}. Memoization falls squarely within this class of refactorings; MemoizeLib is explicitly designed as its drop-in target.

On the measurement side, modern Android performance research has converged on Jetpack Benchmark~\cite{jetpackbenchmark} as the toolchain for both microbenchmarks and macrobenchmarks. It provides warm-up, clock stabilisation, \aot compilation, and outlier rejection, and pairs naturally with \odpm-based energy measurement through Perfetto~\cite{perfetto}. The benchmark harness presented in Section~\ref{sec:study} follows this pattern, including the two-module layout of a library under test and a separate benchmark module with \code{testBuildType~=~release}.

\subsection{Automated detection and refactoring of memoization}

The prior work most directly adjacent to ours is iMMutator~\cite{immutator}, which identifies memoization candidates in \jvm code through a combination of purity analysis and subproblem-overlap heuristics. iMMutator's evaluation demonstrates that the detection step achieves high precision, but that the application step---the emission of correct caching code for each candidate---remains largely unaddressed. This detection-application gap is the explicit motivation for MemoizeLib: by exposing an annotation plus a tested runtime as the refactoring target, a detector such as iMMutator can close the loop from identified candidate to merged patch without owning the correctness of the caching wrapper itself. The same argument applies to \llm-driven performance refactoring~\cite{llm_refactoring,hall_costa}, in which the \llm's correctness surface for memoization is reduced from ``emit a caching wrapper for this class'' to ``select annotation arguments''. In both cases, the library is intended as the trusted target: tested once within the library, rather than regenerated fresh at every call site. The benchmark harness presented in Section~\ref{sec:study} quantifies the extent to which this trust is justified in practice.
